\chapter{Evolução e propostas de continuidade}\label{ap:notas-historicas}
Os registros a seguir preservam a evolução descrita no relatório de referência, com suas datas e fontes. Propostas de integração, relatos de lotação e análises de intenção de viagem não são funcionalidades entregues. As sínteses atuais estão no capítulo de resultados e os limites na Seção~\ref{sec:validade}.

A extensão de São Paulo, examinada na revisão \texttt{4992027}, de 7 de outubro de 2026, acrescentou GTFS/SPTrans e consultas ao Olho Vivo quando havia correspondência de linha e parada. Esse registro histórico não equivale ao fornecimento de dados operacionais de Mauá nem confunde programação com previsão em tempo real.

O fluxo de origem, destino, favoritos e resultados progressivos está documentado na Seção~\ref{sec:fluxo-implementado}. Na revisão histórica, o mapa iniciava no passageiro, com raio mínimo de 250 m no menor eixo e ação de recentralização. Coordenadas simuladas configuráveis substituíam o GPS nos cenários de reprodução; sem elas, utilizava-se a localização autorizada pelo navegador.

O toque prolongado de aproximadamente 0,65 segundo selecionava coordenadas segundo a tarefa ativa (origem ou destino), cancelando a seleção em arrastos ou gestos com dois dedos. O cálculo incluía as caminhadas de acesso. As regras de alteração do embarque, confirmação de cancelamento e preservação de favoritos e enquadramento estão documentadas na Seção~\ref{sec:fluxo-implementado} e no Apêndice~\ref{ap:engenharia}. Fechar o painel apenas mudava a visualização.

A ampulheta sinalizava o cálculo; ao concluir, o recálculo manual ignorava o cache. A busca não reiniciava automaticamente a cada 30 segundos nem a cada atualização do GPS: a atualização periódica pertencia à consulta de próximos ônibus. Nessa consulta, frequência significava intervalo entre veículos e faixa de vigência, com união de faixas consecutivas iguais, não contagem regressiva. A identidade visual própria mantinha a identificação das fontes na documentação e na aplicação.

A distinção entre o demonstrador simulado de Mauá, a implementação funcional, a apreciação profissional de outubro e as verificações posteriores é retomada no capítulo de resultados. O fluxo técnico implementado não comprova redução do esforço do passageiro; essa avaliação permanece pendente. A comparação com Cittamobi, QRCodeAdesivo e Olho Vivo está na fundamentação e não estabelece superioridade do PaP. Mockups anteriores são registros de prototipação, não evidências da solução final.

Os registros recebidos sobre Santos permitiram confrontar o recorte do projeto com uma solução em operação. No texto reproduzido pelo grupo e atribuído à resposta ao protocolo 7657-2026, o Município informa que todos os pontos de parada dispõem de QR Code para consulta das previsões de chegada e que o aplicativo Quanto Tempo Falta oferece linhas, itinerários, localização dos veículos e previsões. A resposta considera desnecessária a implantação de outra sistemática para a finalidade descrita na solicitação, mas reconhece a pertinência de estudos acadêmicos sobre dados e aprimoramento do transporte. Essa posição é registrada como manifestação recebida, sem representar autorização de implantação, parceria ou avaliação do PaP\footnote{Texto atribuído à resposta do Município de Santos ao protocolo 7657-2026, reproduzido pelo grupo. O documento original não integra o material consultado; data e setor remetente não foram informados. As afirmações são relatadas conforme essa reprodução, sem verificação independente da cobertura do serviço.}.

Uma passageira consultada informalmente confirmou utilizar o aplicativo e associou suas previsões à decisão de quando sair de casa. Também relatou situações em que os veículos pareciam mais próximos do que estavam ou deixavam de aparecer no mapa. O relato evidencia utilidade percebida e dificuldades na experiência individual, mas não permite estimar frequência de falhas, determinar suas causas ou generalizar o desempenho do serviço. A resposta afirmativa a uma pergunta que reunia favoritos, rota mais rápida e QR Code não foi tomada como comprovação de todas essas funcionalidades, pois elas não foram verificadas separadamente\footnote{Conversa informal não publicada com uma passageira, relatada pelo grupo a partir de capturas de mensagens. Data não informada; identidade omitida. Não constitui avaliação do PaP.}.

Posteriormente, o grupo informou ter recebido um vídeo no qual a leitura do QR Code abre uma página web do ponto, sem exigir instalação do aplicativo. Segundo a descrição fornecida, o fluxo mostrado não oferece seleção de destino. O registro disponível neste trabalho é essa descrição, e não uma análise independente do arquivo audiovisual. A ausência da opção no percurso descrito não comprova sua inexistência em todas as telas ou versões do serviço. Ainda assim, ela orienta uma questão verificável: como apoiar quem reconhece o destino, mas não sabe quais linhas do ponto podem atendê-lo? A relevância para visitantes e turistas permanece uma hipótese de uso, a ser investigada com participantes desse perfil\footnote{Descrição de vídeo comunicada pelo grupo, sem data informada.}.

Essas informações levaram a uma revisão do argumento do projeto. A entrada pelo QR Code e a consulta sem instalação não são diferenciais exclusivos do PaP. O interesse da proposta passa a ser examinar, no contexto de Mauá, um fluxo que preserve a origem identificada e permita escolher um destino, comparar caminhos e compreender as conexões. Trata-se de uma hipótese de contribuição a avaliar, e não de superioridade demonstrada sobre Santos ou Cittamobi. O aprendizado documentado foi a revisão do escopo a partir da comparação; não houve, até o fechamento desta versão, validação do protótipo com os participantes desses relatos.

Depois de reconhecer a parada, o passageiro ainda precisa relacionar as linhas disponíveis ao lugar a que deseja chegar. A seleção do destino organiza essa etapa da consulta. No cenário São Paulo, essa escolha aciona o planejamento baseado no GTFS e destaca os trechos relevantes. O planejamento admite conexões que não atendem diretamente à origem e considera partidas posteriores ao veículo mais próximo. O catálogo importado permanece consultável sem que sua cobertura parcial seja tomada como rede completa apta ao cálculo.

À luz de Haylett, a interface de São Paulo já apresenta um resumo comparável das alternativas, com duração, sentido e quantidade de trocas, antes dos detalhes de cada trecho. A lista de próximas partidas deve preservar a distinção entre linha, variante e sentido, evitando que o agrupamento visual esconda diferenças de percurso. A verificação de compreensão dessas informações é parte da avaliação proposta, não um benefício já medido.

A arquitetura React/Leaflet, Django/Python e SQLite é apresentada na Figura~\ref{fig:arquitetura-pap}; o índice GTFS, separado dos modelos Django históricos de Mauá, na Figura~\ref{fig:banco-mbm}. A evolução concorrente de 1 de novembro está na Seção~\ref{sec:atualizacao-20261101}. Cobertura e desempenho são sintetizados nas Seções~\ref{sec:cobertura-revisada} e~\ref{sec:desempenho-revisado}; essas medições locais não substituem o protocolo controlado proposto.

No cenário histórico de Mauá, a simulação possui administração separada, com autenticação, para ativação, relógio e velocidade. O passageiro não precisa habilitá-la na tela principal. Quando desligada, veículos e previsões demonstrativas ficam indisponíveis; não há substituição automática por previsões ou posições operacionais de Mauá. O catálogo continua acessível. O isolamento dos cenários permite repetir a demonstração e comparar o comportamento após alterações.

A possibilidade de aproveitar um aplicativo já utilizado na cidade amplia a discussão sobre a continuidade da proposta: o acesso pelo ponto poderia encaminhar o passageiro ao serviço existente, desde que a integração preserve a parada identificada. Em uma implantação futura, o QR Code poderá conter uma URL HTTPS estável associada ao identificador da parada. Esse endereço poderá abrir diretamente a consulta de linhas da parada em um aplicativo compatível já instalado, por mecanismos de \textit{deep linking}, como Android App Links, que associam um domínio verificado ao aplicativo \citep{androidapplinks}. O QR Code identifica o contexto; o sistema operacional e a configuração do aplicativo determinam o tratamento do link.

Quando o aplicativo não estiver instalado, a URL poderá apresentar uma página web com a consulta disponível e uma opção de instalação pela loja. O download e a instalação dependem da ação do usuário. Retomar automaticamente a mesma parada após instalar exige um mecanismo adicional de preservação do contexto, não garantido pelo link comum; uma alternativa simples é ler novamente o QR Code após a instalação. O acesso web permanece disponível no escopo do PaP para que instalar um aplicativo não seja condição para a consulta inicial.

Para abrir especificamente uma parada no Cittamobi, será necessário confirmar com o fornecedor a existência de links suportados, os identificadores de paradas e as condições de integração. Um redirecionamento no servidor, sozinho, não cria uma função que o aplicativo não reconheça. A proposta corresponde a uma entrada contextualizada para o serviço existente; não implica implementação Android nesta etapa nem integração já realizada.

Em deslocamentos recorrentes, informar novamente o mesmo destino pode acrescentar etapas à consulta. Para explorar uma forma de simplificar esse uso, o protótipo permite salvar destinos com nomes como Casa, Trabalho e Escola, com seleção de paradas favoritas na escolha de embarque ou destino. Os dados ficam no navegador, sem conta ou sincronização entre aparelhos. O recurso é opcional: o passageiro ocasional pode iniciar uma consulta pelo QR Code e escolher um destino no mapa sem cadastrar favoritos.

Na implementação de 1 de novembro de 2026, o planejamento combina ônibus, metrô e trens, incluindo Metrô e CPTM. A interface identifica linhas ferroviárias por cores e padrões próprios e representa estações, trechos e transferências entre linhas e modos. A avaliação inclui acessos a pé, espera, circulação e conexões internas estimadas nos complexos cadastrados. Os tempos ferroviários utilizam horários programados ou estimativas de frequência e circulação; o modelo considera paradas, aceleração e frenagem como parâmetros de planejamento, sem equivaler a medição operacional. Permanecem pendentes a validação desses tempos em operação, previsões e posições ferroviárias em tempo real e integração tarifária. Para os ônibus municipais de São Paulo, a consulta de previsões e posições pela API Olho Vivo está implementada, quando o serviço e a correspondência cadastral estão disponíveis. A Seção~\ref{sec:atualizacao-20261101} registra os resultados e limites desta versão.

A conferência ferroviária inicial de 2 de novembro de 2026 confirmou que esse planejamento alcança também o município de Mauá pelas estações Mauá (identificador \texttt{18935}) e Guapituba (\texttt{18937}). Ambas integram as sequências dos dois sentidos da rota \texttt{CPTM L10}, Linha 10--Turquesa, cadastrada como Rio Grande da Serra--Barra Funda. Os dados procedem do pacote GTFS da SPTrans preservado em \nolinkurl{fonts/sptrans_gtfs/} e importado em 6 de outubro de 2026. Os registros de estações, viagens, frequências, calendário e traçados conferem com o banco \nolinkurl{sptrans.sqlite3}; o leitor da materialização compacta também retornou alternativas válidas para esses pontos. Portanto, trata-se de cobertura ferroviária funcional do planejador em Mauá, distinta do demonstrador histórico dos ônibus municipais.

A cobertura temporal depende do cadastro e das regras da consulta. No arquivo examinado, o serviço \texttt{USD} da Linha 10 tem vigência de 1 de outubro de 2023 a 1 de maio de 2027, com os sete dias da semana habilitados e sem exceções cadastradas. As frequências, sem partidas exatas, variam entre 5, 8 e 10 minutos em faixas de 04h00 a 23h59 na origem de cada sentido; a disponibilidade em cada estação depende também do deslocamento temporal na sequência. O código limita novos embarques ferroviários à janela anterior à meia-noite, com abertura às 04h para a Linha 10 e às 04h40 para a Linha 2, e usa horizonte de busca de seis horas. Esses são limites do modelo inspecionado, não horários de operação confirmados em campo. Calendários ativos, partidas alcançáveis e transferências compatíveis continuam necessários; a existência de um par materializado não garante viagem em qualquer horário.

Foi executada uma consulta ao código da API com os bancos reais, em modo somente leitura e sem chamadas externas, entre Mauá e Trianon--Masp (\texttt{18859}), iniciada em 2 de novembro às 00h47min30s, no fuso de São Paulo. A resposta retornou Linha 10 até Tamanduateí e Linha 2 até Trianon--Masp, com primeiro embarque calculado às 04h26 e chegada estimada às 05h32min30s. Em consultas adicionais ao leitor de candidatos, com partida de referência fixada às 08h do mesmo dia, foram obtidas viagens de ida e volta entre Trianon--Masp e Mauá, com duração estimada de 68 minutos, e entre Trianon--Masp e Guapituba, com 75 minutos. Esses horários de referência não são instantes de observação de veículos. Uma consulta de controle após a vigência do calendário não encontrou alternativa temporalmente válida. A revisão de código \texttt{34a46f29}, suas alterações locais, hashes dos fontes, respostas, parâmetros e horários estão preservados em \nolinkurl{auditoria-alcance-ferroviario-2026-11-02/}, junto ao relatório; o identificador Git isolado não reproduz as alterações locais.

O uso a partir de Mauá requer selecionar uma dessas estações como origem ou destino, ou dispor de acesso externo a pé verificado dentro do limite de 1.000 metros do modelo. Esse limite de busca não define uma área de atendimento populacional. O percurso entre a coordenada GTFS e a plataforma não tem tempos de acesso individualmente cadastrados; as transferências internas utilizam estimativas. A consulta ferroviária não depende de placa instalada: pode ser iniciada por seleção no mapa ou link da estação. Não houve implantação física comprovada dos QR Codes, integração operacional com os ônibus municipais de Mauá ou validação operacional dos tempos ferroviários. Também não se inferem cobertura integral do território, percentual de habitantes atendidos ou benefício mensurado para a população a partir dessas verificações técnicas.


Como possibilidade de trabalho futuro, ainda não implementada, propõe-se investigar relatos voluntários dos passageiros sobre a lotação dos veículos. A proposta retoma uma ideia levantada na entrevista inicial, agora considerada à luz da identificação do veículo disponibilizada pela integração com o Olho Vivo. Essa informação seria colaborativa e autodeclarada, não fornecida pela API Olho Vivo, cuja documentação pública não apresenta dados de ocupação. Os relatos poderiam ser associados ao veículo específico pelo prefixo disponibilizado pela API \citep{sptransapi}. Eventuais categorias de lotação e seus critérios dependeriam de definição e avaliação posteriores. A proposta precisaria considerar validade temporal, agregação de múltiplos relatos, qualidade da informação, prevenção de abuso e privacidade, além de avaliar sua utilidade com passageiros.

Um aplicativo Android permanece fora do escopo atual. Sua conveniência será avaliada após a validação do fluxo web, considerando necessidades que não sejam atendidas pelo navegador. A futura aplicação poderá reutilizar a API Django, mas tecnologia, distribuição e eventual sincronização ainda não foram definidas.

A entrevista registrou interesse pelo acesso por QR Code e a sugestão de formalizar a proposta para discussão com o fornecedor. Em resposta a esse encaminhamento, foi preparado um PDF com a sequência dos mockups, explicações do fluxo e perguntas para avaliação. O interlocutor já respondeu por e-mail à pergunta sobre dados de embarque e desembarque. O grupo busca outro interlocutor, diante da ausência de resposta aos contatos recentes, e ainda não realizou avaliação de uso documentada junto à comunidade de Mauá; a resposta técnica sobre bilhetagem não é tratada como teste ou aprovação dos mockups. A análise dessa segunda interação será registrada separadamente da entrevista de definição do problema  (entrevista com Osmar Junior, 1135--1163 s; 2293--2425 s).

O protótipo web funcional é o principal material para apresentação e avaliação da solução atual. O PDF de mockups e o Figma documentam a etapa de prototipação e a evolução da interface. A apresentação deverá identificar a revisão da aplicação examinada, o cenário e as informações disponíveis nessa versão. A discussão priorizará a contribuição do QR Code para identificar a parada em comparação com o Cittamobi, a compreensão do fluxo, a viabilidade percebida e os limites da eventual integração. O retorno profissional do interlocutor será separado dos resultados de tarefas realizadas por passageiros.

Também deverão ser levantados os responsáveis e as condições de acesso ao cadastro operacional: identificadores e coordenadas de plataformas por sentido; linhas e variantes vigentes; sequência de paradas e geometrias; horários, calendários e exceções; regras de transferência; localização e previsão, quando disponíveis; periodicidade de atualização e condições de uso. Para a totalidade dos ônibus municipais de Mauá, o acesso operacional a esses dados ainda não está confirmado; a ampliação metropolitana ARTESP está descrita na Seção~\ref{sec:artesp-maua}. Em São Paulo, o protótipo já consulta o cadastro GTFS e a API Olho Vivo, sujeitos à cobertura e disponibilidade dos serviços. Qualquer piloto físico exigirá definição de responsáveis, autorização para as placas e manutenção dos QR Codes e das informações.

A discussão com a comunidade também suscitou uma questão para além da consulta individual: de que modo os destinos informados poderiam contribuir para compreender necessidades de deslocamento? Para explorar essa possibilidade, o grupo perguntou quais informações de embarque e desembarque já eram registradas. Em resposta ao e-mail, o interlocutor esclareceu que, no contexto operacional descrito, a bilhetagem registra a validação do cartão, sua data e hora e a linha ou o veículo, podendo associar uma posição aproximada ao embarque quando integrada ao GPS. Segundo o interlocutor, o desembarque não é registrado diretamente, e, no Terminal Central de Mauá, a validação ocorre no acesso à área paga, sem identificar necessariamente o veículo utilizado depois. Essas informações constituem relato do interlocutor sobre o sistema discutido, não uma conclusão sobre todos os sistemas de bilhetagem\footnote{E-mail não publicado de Osmar Junior a André Meneghelli Vale, enviado em 2 set. 2026, às 19h42, conforme cabeçalho da captura fornecida pelo destinatário. Resposta à consulta sobre bilhetagem e informações de origem e destino.}.

O retorno também menciona pesquisas de ``sobe e desce'' e inferências a partir de embarques em viagens posteriores como meios de apoiar o planejamento. A contagem de embarques e desembarques por parada descreve movimentos agregados, mas não associa necessariamente cada origem ao destino de uma mesma pessoa. O e-mail de 2 de setembro de 2026, às 19h42, conforme o cabeçalho da captura fornecida pelo grupo, reforça a pertinência de investigar uma fonte complementar de informação: os pares de origem e destino informados durante o planejamento no PaP (e-mail de Osmar Junior, 2 set. 2026).

Ao ler o QR Code, a consulta recebe o identificador da parada de origem; ao escolher um destino, o usuário informa uma intenção de deslocamento. O agrupamento dessas consultas por origem, destino e faixa horária poderia indicar destinos procurados, relações recorrentes entre regiões e necessidades de conexão. Essa contribuição potencial amplia o interesse do acesso contextualizado: além de facilitar a consulta, ele pode produzir dados para investigar a demanda informacional dos usuários.

Esses registros deverão ser denominados \textit{consultas ou intenções de viagem}, e não viagens realizadas. Abrir o QR Code não comprova embarque; selecionar um destino não comprova desembarque. Uma pessoa pode comparar vários destinos, consultar repetidamente ou desistir da viagem, e uma consulta pode atender a mais de um passageiro. O conjunto representa apenas quem utiliza a ferramenta, sujeito a diferenças de acesso à internet e de adesão. Portanto, contagens de consultas não equivalem a contagens de passageiros nem a uma matriz origem-destino representativa da população.

Antes de coletar dados para essa finalidade, deverão ser definidos objetivo, informação ao usuário, participação voluntária, retenção e acesso aos registros. Uma estrutura mínima a avaliar incluiria identificadores da origem e do destino, faixa de tempo e identificação de consultas de teste ou demonstração, separadas das demais consultas de uso, sem exigir nome, cartão de transporte ou rastreamento contínuo. Destinos salvos como Casa e Trabalho não deverão ser publicados individualmente; mesmo sem nomes, pares precisos de locais podem permitir associações pessoais. A divulgação deverá priorizar agrupamentos por região e período, com tratamento de grupos pouco frequentes e regras de qualidade para repetições e consultas de teste.
